\section{Page Fault 集成}
\label{sec:vma-pf}

\#PF handler（ISR 14）收到故障地址 \texttt{cr2} 和错误码 \texttt{err\_code} 后，
调用 \texttt{vma\_find(cr2)} 可以产生三种诊断结果：

\begin{figure}[H]
  \centering
  % figures/ch11/fig10-pf-vma-diag.tex — auto-extracted from chapter source
\begin{tikzpicture}[
  box/.style={draw, rounded corners=3pt, minimum height=0.7cm,
              minimum width=4.5cm, font=\small, align=center},
  decision/.style={draw, diamond, aspect=2.5, font=\small,
                   inner sep=2pt, fill=yellow!10},
  arr/.style={-{Stealth[length=5pt]}, thick},
]
  \node[decision] (find) at (0, 0) {\texttt{vma\_find(cr2)}};

  \node[box, fill=red!12]    (no)   at (-4.5, -2) {Case 1: 返回 NULL\\非法访问 → panic / kill};
  \node[decision]             (perm) at (2, -2) {权限匹配？};
  \node[box, fill=orange!12] (deny) at (0, -4) {Case 2: 权限不匹配\\如写只读区 → SIGSEGV};
  \node[box, fill=green!12]  (ok)   at (5, -4) {Case 3: 合法缺页\\demand-paging → 分配物理页};

  \draw[arr] (find) -- node[left, font=\scriptsize] {NULL} (no);
  \draw[arr] (find) -- node[right, font=\scriptsize] {非 NULL} (perm);
  \draw[arr] (perm) -- node[left, font=\scriptsize] {否} (deny);
  \draw[arr] (perm) -- node[right, font=\scriptsize] {是} (ok);
\end{tikzpicture}

  \caption{\#PF handler 中的 VMA 三种诊断}
  \label{fig:pf-vma-diag}
\end{figure}

\begin{enumerate}
  \item \textbf{Case 1: 无 VMA}（\texttt{vma\_find} 返回 NULL）——
        该地址不属于任何已注册区域，是真正的非法访问。
        内核态触发则 panic；用户态触发则杀死进程。

  \item \textbf{Case 2: 权限不匹配}——地址属于某个 VMA，但操作违反权限。
        例如写入一个仅 \texttt{VMA\_READ} 的区域。
        对应 Linux 的 \texttt{SIGSEGV}。

  \item \textbf{Case 3: 合法缺页（demand paging candidate）}——
        地址属于 VMA 且权限匹配，只是物理页尚未分配。
        分配物理页 → 建立映射 → 返回，程序继续运行。
\end{enumerate}

当前实现中我们尚未启用 demand paging，所有映射在 \texttt{vmm\_alloc\_pages} 时
一步到位。但 VMA 框架为将来的按需分页奠定了基础。

% ============================================================================

